% ==============================================================================
% ch01_overview_principles.tex — บทที่ 1
% ==============================================================================
\chapter{วิวัฒนาการและทฤษฎีการออกแบบสื่อการสอนสามมิติเสมือนจริง}
\label{chap:ch01}

\begin{teachingplan}{แผนการบริหารการสอนประจำบทที่ 1}
\noindent\textbf{หัวข้อเนื้อหาประจำบท}
\begin{enumerate}[leftmargin=1.5em, itemsep=2pt]
    \item ข้อจำกัดของการทดลองทางกายภาพในห้องปฏิบัติการและสื่อสองมิติแบบดั้งเดิม
    \item วิวัฒนาการจาก 2D Canvas สู่ 3D WebGL Virtual Lab และ Augmented Reality (AR)
    \item ทฤษฎีการเรียนรู้เชิงสร้างสรรค์นิยม (Constructivism) และ Active Learning Cycle
    \item ทฤษฎีภาระการรู้คิด (Cognitive Load Theory) กับการออกแบบหน้าจอห้องปฏิบัติการเสมือนจริง
    \item สถาปัตยกรรมระบบ RBRU Next Gen Lab 2026 และมาตรฐานโอลิมปิกวิชาการ IJSO
\end{enumerate}

\noindent\textbf{วัตถุประสงค์เชิงพฤติกรรม (Behavioral Objectives)}
\begin{enumerate}[leftmargin=1.5em, itemsep=2pt]
    \item ผู้เรียนสามารถอธิบายความแตกต่างระหว่างสื่อการสอน 2D, 3D WebGL และ AR ได้อย่างถูกต้อง
    \item ผู้เรียนสามารถวิเคราะห์ข้อจำกัดในการจัดการเรียนการสอนวิทยาศาสตร์แบบเดิม และระบุจุดเด่นของสื่อเสมือนจริงได้
    \item ผู้เรียนสามารถประยุกต์ทฤษฎี Cognitive Load ในการจัดวางส่วนติดต่อผู้ใช้ (UI/UX) ของห้องปฏิบัติการเสมือนจริงได้
    \item ผู้เรียนสามารถเขียนผังมโนทัศน์ (Conceptual Framework) ในการพัฒนาสื่อการสอน 3 มิติของตนเองได้
\end{enumerate}

\noindent\textbf{กิจกรรมการเรียนการสอน (Teaching-Learning Activities)}
\begin{enumerate}[leftmargin=1.5em, itemsep=2pt]
    \item บรรยายนำเข้าสู่บทเรียนด้วยกรณีศึกษาการทดลองที่อันตรายหรือไม่สามารถมองเห็นด้วยตาเปล่า (30 นาที)
    \item สาธิตการใช้งานระบบ RBRU Next Gen Lab 2026 บนเบราว์เซอร์และการจำลองวงจรหกเหลี่ยม (30 นาที)
    \item กิจกรรมกลุ่มย่อย: วิเคราะห์ตัวชี้วัดการเรียนรู้ฟิสิกส์เพื่อคัดเลือกหัวข้อสร้างสื่อเสมือนจริง (30 นาที)
\end{enumerate}

\noindent\textbf{สื่อการเรียนการสอน}
\begin{enumerate}[leftmargin=1.5em, itemsep=2pt]
    \item สไลด์บรรยายอิเล็กทรอนิกส์, เอกสารคู่มือบทที่ 1, เว็บไซต์สาธิตระบบ Next Gen Lab
\end{enumerate}

\noindent\textbf{การวัดและประเมินผล}
\begin{enumerate}[leftmargin=1.5em, itemsep=2pt]
    \item ประเมินจากการมีส่วนร่วมในกิจกรรมกลุ่มย่อยและการนำเสนอแนวคิดสื่อการสอน
    \item ประเมินจากการตอบคำถามทบทวนท้ายบทที่ 1 ได้คะแนนไม่น้อยกว่าร้อยละ 80
\end{enumerate}
\end{teachingplan}

\newpage

\section{บทนำและข้อจำกัดของห้องปฏิบัติการกายภาพ}
ในการจัดการเรียนรู้วิทยาศาสตร์ระดับมัธยมศึกษาและอุดมศึกษา “การทดลองในห้องปฏิบัติการ” ถือเป็นหัวใจสำคัญในการสร้างเสริมทักษะกระบวนการทางวิทยาศาสตร์ ทว่าในทางปฏิบัติ สถานศึกษาจำนวนมากต้องเผชิญกับอุปสรรคสำคัญหลายประการ:
\begin{itemize}[leftmargin=1.5em, itemsep=3pt]
    \item \textbf{ความเสี่ยงด้านความปลอดภัยและอุปกรณ์ชำรุด:} การทดลองที่เกี่ยวข้องกับไฟฟ้าแรงสูง สารเคมีอันตราย หรือเครื่องมือวัดความแม่นยำสูง (Precision Galvanometers) มักเกิดความเสียหายได้ง่ายจากการต่อวงจรผิดพลาดของนักเรียน
    \item \textbf{ข้อจำกัดด้านงบประมาณและเวลา:} งบประมาณในการจัดซื้อชุดทดลองราคาแพงไม่เพียงพอต่อสัดส่วนนักเรียน ส่งผลให้นักเรียนส่วนใหญ่ทำได้เพียงเฝ้ามองตัวแทนกลุ่มทดลองเท่านั้น
    \item \textbf{ข้อจำกัดของสิ่งที่ไม่สามารถมองเห็นได้ด้วยตาเปล่า (Unobservable Phenomena):} การไหลของอิเล็กตรอนในวงจร, การเหนี่ยวนำแม่เหล็กไฟฟ้า, หรือสนามศักย์ไฟฟ้า ล้วนเป็นนามธรรมที่ภาพวาด 2 มิติในหนังสือเรียนไม่สามารถถ่ายทอดพลวัตได้อย่างแท้จริง
\end{itemize}

\begin{figure}[h!]
\centering
\begin{tikzpicture}[scale=0.9, every node/.style={scale=0.9}]
    % 2D Flat Plane Box
    \draw[fill=gray!15, draw=gray!60, rounded corners=6pt] (0,0) rectangle (4.2,3.5);
    \node at (2.1,3.1) {\bfseries\color{rbruNavy} สื่อ 2D แบนราบแบบดั้งเดิม};
    \draw[thick, red] (1,1) rectangle (3.2,2.2);
    \node at (2.1,1.6) {\small ภาพนิ่งในหนังสือ};
    \node[align=center, text width=3.8cm, font=\footnotesize, color=gray!80!black] at (2.1,0.5) {มองได้มุมเดียว ขาดมิติความลึก\\ไม่สามารถโต้ตอบได้};

    % Arrow 1
    \draw[-{Stealth[scale=1.2]}, thick, rbruBlue] (4.5,1.75) -- (5.8,1.75) node[midway, above, font=\scriptsize] {Web 3D};

    % 3D Virtual Lab Box
    \draw[fill=blue!5, draw=rbruBlue, rounded corners=6pt] (6.0,0) rectangle (10.2,3.5);
    \node at (8.1,3.1) {\bfseries\color{rbruBlue} 3D Virtual Lab};
    % 3D isometric cube
    \draw[thick, fill=cyan!20, draw=cyan!80!black] (7.2,1.2) -- (8.6,1.2) -- (8.6,2.2) -- (7.2,2.2) -- cycle;
    \draw[thick, fill=cyan!35, draw=cyan!80!black] (7.2,2.2) -- (7.8,2.7) -- (9.2,2.7) -- (8.6,2.2) -- cycle;
    \draw[thick, fill=cyan!10, draw=cyan!80!black] (8.6,1.2) -- (9.2,1.7) -- (9.2,2.7) -- (8.6,2.2) -- cycle;
    \node[align=center, text width=3.8cm, font=\footnotesize, color=rbruNavy] at (8.1,0.5) {หมุนกล้อง 360 องศา\\จำลองฟิสิกส์สมจริง};

    % Arrow 2
    \draw[-{Stealth[scale=1.2]}, thick, rbruEmerald] (10.4,1.75) -- (11.7,1.75) node[midway, above, font=\scriptsize] {Spatial AI};

    % AR MediaPipe Box
    \draw[fill=green!5, draw=rbruEmerald, rounded corners=6pt] (11.9,0) rectangle (16.1,3.5);
    \node at (14.0,3.1) {\bfseries\color{rbruEmerald} Touchless AR Lab};
    \draw[thick, rbruEmerald] (13.2,1.3) circle (0.5cm);
    \draw[thick, rbruGold] (14.8,1.3) circle (0.5cm);
    \draw[thick, dashed, cyan] (13.2,1.8) -- (14.8,1.8);
    \node at (14.0,2.1) {\tiny Pinch Gesture};
    \node[align=center, text width=3.8cm, font=\footnotesize, color=rbruNavy] at (14.0,0.5) {สั่งการด้วยมือเปล่าผ่านกล้อง\\ประสานโลกจริงกับโลกเสมือน};
\end{tikzpicture}
\caption{วิวัฒนาการของการจัดการเรียนรู้จากสื่อสองมิติ สู่ 3D WebGL และ Augmented Reality (AR)}
\label{fig:evolution_media}
\end{figure}

\section{ทฤษฎีการเรียนรู้และภาระการรู้คิด (Cognitive Load Theory)}
การพัฒนาสื่อเสมือนจริงที่มีประสิทธิภาพ มิใช่เพียงการสร้างภาพกราฟิกที่สวยงามตระการตาเท่านั้น หากแต่ต้องยึดหลักจิตวิทยาการเรียนรู้ เพื่อให้สมองของผู้เรียนสามารถประมวลผลข้อมูลได้อย่างสูงสุดตามทฤษฎี **Cognitive Load Theory (Sweller, 1988)**:

\begin{theoremlab}{ภาระการรู้คิด 3 มิติในการออกแบบ Virtual Laboratory}
ภาระการรู้คิดรวมของมนุษย์ ($CL_{\text{Total}}$) ประกอบด้วย 3 องค์ประกอบหลัก:
\begin{equation}
CL_{\text{Total}} = CL_{\text{Intrinsic}} + CL_{\text{Extraneous}} + CL_{\text{Germane}}
\end{equation}
\begin{itemize}[leftmargin=1.5em, itemsep=2pt]
    \item \textbf{Intrinsic Load ($CL_{\text{Intrinsic}}$):} ความยากตามธรรมชาติของเนื้อหา (เช่น สูตรการยุบวงจร $R_{\text{ชุด}} = \frac{r_1 r_2}{r_1+r_2}$) ผู้สอนต้องซอยย่อยเนื้อหาเป็นขั้นตอน (Scaffolding Steps 1--4)
    \item \textbf{Extraneous Load ($CL_{\text{Extraneous}}$):} ภาระส่วนเกินที่เกิดจากส่วนต่อประสานที่สับสน แผงควบคุมบดบังอุปกรณ์ หรือปุ่มกดยาก \textbf{ต้องลดทอนให้เหลือน้อยที่สุด} โดยใช้ดีไซน์ Glassmorphism โปร่งแสง และแผงควบคุมที่พับเก็บ/ย่อขยายได้อิสระ
    \item \textbf{Germane Load ($CL_{\text{Germane}}$):} การจัดสรรพลังสมองเพื่อสร้างโครงสร้างความรู้ใหม่ (Schema Construction) สื่อต้องเสริมมิติภาพเคลื่อนไหว (อนุภาคกระแสวิ่ง, เข็มมิเตอร์แกว่ง) เพื่อกระตุ้นมโนทัศน์ที่ถูกต้อง
\end{itemize}
\end{theoremlab}

\begin{tipbox}{กฎเหล็กการป้องกันการบดบังทางสายตา (Visual Occlusion Rule)}
ในสื่อ 3 มิติเสมือนจริง อย่าให้ข้อความอธิบายทับซ้อนกับอุปกรณ์ทดลองเด็ดขาด ในระบบ RBRU Next Gen Lab เราออกแบบให้แถบด้านข้าง (Sidebars) ใช้คุณสมบัติ \texttt{backdrop-filter: blur(14px)} โปร่งแสง และมีปุ่มพับเก็บ (Collapse) เพื่อให้ผู้เรียนมองเห็นพฤติกรรมของอุปกรณ์ 3D ที่พื้นหลังได้ตลอดเวลา
\end{tipbox}

\section{สถาปัตยกรรมระบบ RBRU Next Gen Lab 2026 และ 7 กลุ่มวิชาฟิสิกส์}
ระบบห้องปฏิบัติการเสมือนจริง RBRU Next Gen Lab 2026 ได้รับการออกแบบให้ทำงานในรูปแบบ **Standalone Client-Side Web Application** กล่าวคือ รันได้ทันทีบนเว็บเบราว์เซอร์มาตรฐานสากล (Chrome, Safari, Edge, Firefox) โดยไม่ต้องติดตั้งโปรแกรมเสริม และได้รับการขยายขอบเขตการทดลองครอบคลุม **7 กลุ่มวิชาฟิสิกส์พื้นฐาน (7 Fundamental Physics Domains)** อย่างเป็นระบบ:

\begin{enumerate}[leftmargin=1.5em, itemsep=3pt]
    \item \textbf{ไฟฟ้าและวงจรไฟฟ้า (Electricity \& Circuits):} ศึกษาทฤษฎีวงจรไฟฟ้ากระแสตรง กฎของโอห์ม กฎของเคอร์ชอฟฟ์ วงจรลดรูป 6 เหลี่ยม วงจรบริดจ์ และวงจรตัวเก็บประจุ RC
    \item \textbf{แม่เหล็กไฟฟ้าและการเหนี่ยวนำ (Electromagnetism \& Induction):} ศึกษาการเหนี่ยวนำแม่เหล็กไฟฟ้า กฎของฟาราเดย์และเลนซ์ สนามแม่เหล็กในขดลวดโซเลนอยด์ และขดลวดเฮล์มโฮลทซ์
    \item \textbf{กลศาสตร์และการเคลื่อนที่ (Mechanics \& Dynamics):} ศึกษาการเคลื่อนที่แบบฮาร์มอนิกอย่างง่าย ลูกตุ้มเพนดูลัม รางลมไร้แรงเสียดทาน และการตกอิสระวัดค่าความเร่ง $g$
    \item \textbf{คลื่นและเสียง (Waves \& Acoustics):} ศึกษาปรากฏการณ์การสั่นพ้องในหลอดเรโซแนนซ์ คลื่นนิ่งในเส้นเชือก และความเร็วเสียงในอากาศ
    \item \textbf{ทัศนศาสตร์และแสง (Optics \& Wave Optics):} ศึกษาการเกิดภาพบนรางทัศนศาสตร์ กฎของสเนลล์ การแทรกสอด และการเลี้ยวเบนของแสงผ่านเกรตติง
    \item \textbf{อุณหพลศาสตร์และสสาร (Thermodynamics \& Matter):} ศึกษากฎของก๊าซอุดมคติ แคลอริมิเตอร์หาความจุความร้อนจำเพาะ และการขยายตัวเชิงเส้นของโลหะ
    \item \textbf{ฟิสิกส์ยุคใหม่และควอนตัม (Modern \& Quantum Physics):} ศึกษาปรากฏการณ์โฟโตอิเล็กทริก ครึ่งชีวิตการสลายตัวของกัมมันตรังสี และการทดลองของฟรองค์-เฮิร์ตซ์
\end{enumerate}

\begin{activitybox}{การสำรวจโครงสร้างโค้ดและการเปิดใช้งานระบบ}
เปิดดูโฟลเดอร์โครงงาน \texttt{nextgen\_lab2026/} จะพบว่าโครงสร้างไฟล์ถูกแบ่งออกเป็นโมดูลอิสระ:
\begin{itemize}[leftmargin=1.5em, itemsep=2pt]
    \item \texttt{index.html}: จุดเริ่มต้นของระบบ โครงสร้าง Semantic และแผง UI
    \item \texttt{css/}: โมดูลสไตล์ชีตแยกย่อย 7 ไฟล์ (\texttt{main}, \texttt{panels}, \texttt{meters}, \texttt{ar}, \texttt{modal}, \texttt{portal}, \texttt{ai-bot})
    \item \texttt{js/}: ตรรกะฟิสิกส์และการเรนเดอร์แยกโมดูลตามการทดลอง และทะเบียน \texttt{experiments-registry.js}
\end{itemize}
ทดลองเปิดไฟล์ \texttt{index.html} ผ่านเบราว์เซอร์ แล้วทดสอบสลับโหมดมิเตอร์ และคลิกเปิดแผง \textbf{Mr. Deep\_xAI Co-creation Assistant}
\end{activitybox}

\section*{คำถามทบทวนและใบงานปฏิบัติการท้ายบทที่ 1}
\addcontentsline{toc}{section}{คำถามทบทวนและใบงานปฏิบัติการท้ายบทที่ 1}

\begin{enumerate}[leftmargin=1.5em, itemsep=8pt]
    \item จงอธิบายว่าการจำลองการทดลองด้วย 3D WebGL สามารถช่วยแก้ไขข้อจำกัดของห้องปฏิบัติการวิทยาศาสตร์กายภาพในโรงเรียนได้อย่างไร ยกตัวอย่างปรากฏการณ์ที่มองไม่เห็นด้วยตาเปล่ามาอย่างน้อย 2 ตัวอย่าง
    \item ตามทฤษฎี Cognitive Load Theory หากผู้ออกแบบสร้างหน้าจอที่มีปุ่มและกราฟิกกระพริบรกทึบ จะส่งผลกระทบต่อภาระการรู้คิดตัวใด ($CL_{\text{Intrinsic}}, CL_{\text{Extraneous}}$ หรือ $CL_{\text{Germane}}$) และมีแนวทางแก้ไขในเชิงสถาปัตยกรรมซอฟต์แวร์อย่างไร
    \item \textbf{ใบงานปฏิบัติการ 1.1:} ให้นักศึกษาเลือกหัวข้อการทดลองวิทยาศาสตร์ในระดับมัธยมศึกษาที่ตนเองสนใจ 1 หัวข้อ (เช่น การหักเหของแสง, การเคลื่อนที่แบบโพรเจกไทล์, หรือคลื่นนิ่งในเส้นเชือก) แล้วร่างแผนผังส่วนติดต่อผู้ใช้ (UI Wireframe) แบ่งเป็น แผงควบคุมฝั่งซ้าย, ฉากทดลอง 3D กึ่งกลาง, และแผงวิเคราะห์ผลฝั่งขวา พร้อมอธิบายเหตุผลในการจัดวาง
\end{enumerate}
